%% raw_s2.tex -- raw transcription of Slide Set 2 (Linear models; feature engineering; EDA), PDF pages 1-78 (compressed)
\section{Slide Set 2 --- Linear models, feature engineering, EDA}
\ixh{linear model}\ixh{linear regression}\ixh{feature engineering}\ixh{exploratory data analysis (EDA)}
{\color{figgray}\itshape [In Part I: S2.1, 3--5, 7--10, 12--14, 16--17, 20--24, 26--31, 33--36, 38--44, 46--49, 51--52, 54, 56--78.]}

\slide{2}{2}{Required reading}
Undergrads: ch.~2 of \ixc{ISLR}; grads: ch.~2 of \ix{ESL}.

\slide{2}{6}{[Missing-data patterns and mechanisms]}
\fig{Patterns over $Y_1..Y_4$: (a) \emph{univariate}\ixh{missingness pattern}: only $Y_4$ has missing rows; (b) \emph{monotone}: $Y_3$ missing bottom third, $Y_4$ bottom half (each later variable missing on a superset of rows); (c) \emph{general}: scattered gaps everywhere; (d) \emph{planned}: $Y_2,Y_3,Y_4$ missing in staggered blocks by design.}
\fig{Causal DAGs (McElreath), all with A$\to$H, H$\to$H*, D$\to$H* (D = dog = missingness mechanism): MCAR\ixh{missing completely at random (MCAR)}: nothing $\to$ D; MAR\ixh{missing at random (MAR)}: A$\to$D; MNAR\ixh{missing not at random (MNAR)}: H$\to$D.}

\subsection{Linear regression}

\slide{2}{11}{Linear regression}
PSTAT 126. Gauss\ixh{Gauss, Carl Friedrich}, 1795. Simple, interpretable\ixh{interpretability}, often very useful; many nonlinear methods generalize \ix{linear regression}. (1809): method of \ix{least squares}, MLE\ixh{maximum likelihood}, Gaussian distribution.

\slide{2}{15--18}{Simple linear regression; which line?}
\ixh{simple linear regression}
\fig{9 counties, confirmed (50k--200k) vs.\ population (1M--3.3M): $\approx$(0.89M,73k), (0.98M,70k), (1.52M,70k), (1.66M,56k), (1.93M,75k), (2.15M,212k), (2.41M,200k), (3.17M,182k), (3.31M,171k); line $y=\beta_0+\beta_1x$ from $\approx$71k to $\approx$210k. Riverside ($x\approx$2.41M): observed $\approx$200k vs.\ predicted $\approx$157k; gap = \ix{residual}.}

\slide{2}{19}{[Workflow]}
Data (training set\ixh{training data}{} $\{(x_i,y_i)\}_{i=1}^n$) $\to$ Learn ($\hat\beta_0,\hat\beta_1$\ixh[beta hat]{$\hat\beta$}) $\to$ Evaluate (1.~accuracy of $\hat\beta_0,\hat\beta_1$; 2.~accuracy of the model).

\slide{2}{25}{Accuracy of the coefficient estimates}
\ixh{sampling variability}
\figref{10}

\subsection{Multiple linear regression}
\ixh{multiple linear regression}

\slide{2}{32}{Coefficient estimates}
$i$th \ix{residual}: $y_i - \hat y_i = y_i - \hat\beta_0 - \sum_{j=1}^p \hat\beta_j x_{ij}$.
\fig{ISLR 3.4: least-squares \emph{plane}\ixh{regression plane}{} in $(X_1,X_2,Y)$; vertical segments = residuals (points above and below).}

\slide{2}{37}{[Stepwise selection, 5 variables]}
\ixh{stepwise selection}\ixh{forward selection}\ixh{backward selection}
Forward: null $\to$ add most significant ($X_2$) $\to$ add $X_5$ $\to$ \ldots\ until stopping rule or out of variables. Backward: full $\{X_1..X_5\}$ $\to$ remove least significant ($X_4$) $\to$ remove $X_1$ $\to\{X_2,X_3,X_5\}$.

\slide{2}{45}{tidymodels: traditional dummy variables}
\ixch{tidymodels}\ixh{dummy variable}\ixh{categorical predictor}
\ixc{recipes}: ``\ixc{step\_dummy()} creates a \emph{specification} of a recipe step that will convert nominal data (e.g.\ character or factors) into one or more numeric binary model terms for the levels of the original data.''

\subsection{Interactions, polynomial regression}

\slide{2}{50}{Interactions vs.\ main effects}
\ixh{main effect}\ixh{interaction term}
\figref{11}

\slide{2}{53}{Interaction: Credit example}
\ixh{Credit (dataset)}
\figref{12}

\slide{2}{55}{Non-linear relationships}
\ixh{nonlinear relationship}\ixh{polynomial regression}\ixh{overfitting}
\figref{13}
